% Supplementary Notes for the LeafletFA Genome Biology manuscript.
% Same LaTeX conventions as manuscript/methods/methods_CURRENT.tex (\subsection*, tabular, \cite keys).
% Supplementary Note 1 is derived from manuscript/NOTE_duplicate_reads.md (2026-08-15), updated 2026-10-27 to
% relate its estimate to the data-measured median (SpliceVI2026, Supplementary Fig. S1).

\section*{Supplementary Notes}

\subsection*{Supplementary Note 1: Duplicate reads and the effective number of molecules per event}
\label{suppnote:duplicates}

\paragraph{Method.} Smart-seq2 libraries carry no unique molecular identifiers, so the reads a cell
contributes to an ATSE are copies of a small number of cDNA molecules rather than independent
observations. To measure this without fitting any model, we restricted to ATSEs whose pooled composition
within one cell type (microglia, 3{,}000 cells) is clearly mixed, with the most-used junction's pooled
usage between 0.2 and 0.8, so that alternative splicing occurs at the population level. For each cell
covering such an event we asked whether all of its reads fall on a single junction, and compared the
observed fraction with the fraction expected if reads were independent draws from the pooled composition.
Treating the cell's reads as $M$ independent molecules, $P(\text{all reads on one junction}) = p^M + (1-p)^M$
was solved for $M$ at $p \approx 0.5$.

\paragraph{Result.} Across all mixed (cell, ATSE) pairs, 68\% place every read on one junction, at a median of
48 reads per pair; under independent reads this would essentially never happen above 50 reads
(Supplementary Note Table 1). The implied number of independent molecules rises from 1.7 to 3.1 per event
as read counts rise from about 6 to about 124: reads increase 19-fold while molecules increase 2.5-fold,
the signature of PCR duplication. The effect is the regime the whole atlas sits in, not a property of a
few poor cells.

\begin{table}[h]
\centering
\small
\begin{tabular}{rccc}
\hline
\textbf{Reads per cell and event} & \textbf{All reads on one junction} & \textbf{Expected if independent} & \textbf{Effective molecules $M$} \\
\hline
5--8 ($\sim$6)       & 0.632 & 0.054   & 1.66 \\
10--18 ($\sim$14)    & 0.532 & 0.005   & 1.91 \\
20--48 ($\sim$34)    & 0.434 & 0.0001  & 2.20 \\
50--198 ($\sim$124)  & 0.238 & $\approx$0 & 3.07 \\
\hline
\end{tabular}
\caption*{\textbf{Supplementary Note Table 1.} Fraction of cells placing all reads on one junction, in ATSEs
with mixed pooled usage (microglia), by read depth; the fraction expected under independent reads; and the
implied number of independent molecules ($p \approx 0.5$).}
\end{table}

\paragraph{Relation to the data-measured median and to the model's floor.} These estimates come from deeply
covered events with mixed composition, chosen so that duplication is visible. Over all observed events,
most of which have few reads and a skewed composition, the median number of effective molecules measured
from the within-event correlation of reads is lower, about 1.2 (1.19 in Tabula Muris Senis and 1.26 in Allen
Brain) \cite{SpliceVI2026}. The two estimates use different event sets and estimators, and agree that a cell
contributes only one to three independent molecules to an event, far fewer than its read count. LeafletFA's
per-event concentration floor of 2 matches the model-implied value (2.15 and 2.17 in the same figure); it
permits less overdispersion than the data-measured median shows and is a modeling choice, not a
measurement (Methods).

\paragraph{Limits.} The molecule estimate assumes $p \approx 0.5$ on the events used, so it fixes an order of
magnitude rather than a decimal. The all-or-nothing signature does not by itself separate technical
duplication from genuine single-cell isoform commitment; under either, a low concentration is the correct
description of the per-cell evidence.
